\section{The Adjudication Discipline}
\label{sec:discipline}

None of the five rules is novel in isolation; applied together,
\emph{before} any model verdict is recorded, they form a lens-failure
detector with a measured record. Two clauses were codified
mid-catalogue by the incidents that exposed them (Rule 3's
positive-control clause, \S\ref{sec:wronglayer}; Rule 2's graded
crosscheck, Appendix~\ref{app:threelenses}); all five rules were in
force before the decisive adjudication of \S\ref{sec:gauge}.

\textbf{Rule 1: instrument-health adjudication precedes model
verdicts.} Before a sweep's pre-registered endpoint may be read, a
health checklist walks the raw per-cell artifacts: gate violations,
convergence, a contradiction scan between independent readouts (the
rule that fired in \S\ref{sec:gauge}).

\textbf{Rule 2: pre-register a crosscheck that removes the primary
lens's strongest assumption}, graded wherever that assumption provably
fails, the switching rule fixed in advance.

\textbf{Rule 3: falsifier teeth, run to completion.} Every negative
control is executed, not merely written; guards are mutation-tested;
every probe rung carries a planted-signal positive control; the
adjudication itself carries a pre-registered falsifier whose firing
voids it.

\textbf{Rule 4: contradictions are resolved against raw artifacts.}
When two reads disagree, neither recency nor authority decides: both
are recomputed from the committed raw files and the tiebreak is
recorded before anything downstream runs.

\textbf{Rule 5: argmax closure.} Wherever a claim requires exact
continuous recovery, the grading path must not be satisfiable by
argmax decoding, under which a rank-1 matrix can recover on the order
of $d$ associations \citep{nichani2025factual}; the converse gap was
pre-registered before it fired (\S\ref{sec:wronglayer}).
% <!-- evidence: N1 -->
